% 小行星卫星（综述）
% license CCBYSA3
% type Wiki

本文根据 CC-BY-SA 协议转载翻译自维基百科\href{https://en.wikipedia.org/wiki/Minor-planet_moon}{相关文章}。

\begin{figure}[ht]
\centering
\includegraphics[width=6cm]{./figures/fe309a0bf64922bd.png}
\caption{} \label{fig_Minor_2}
\end{figure}
一个小行星卫星是指作为天然卫星环绕一颗小行星运行的天文对象。截至 2022 年 1 月，已有 457 颗已知或被怀疑拥有卫星的小行星。之所以说对小行星卫星（以及一般意义上的双天体系统）的发现十分重要，是因为通过确定它们的轨道，可以估计主星的质量与密度，从而获得通常难以通过其他方式获取的物理性质信息。

在这些卫星中，有若干在相对尺寸上非常巨大：90 号 Antiope、Mors–Somnus 与 Sila–Nunam 的尺寸约为其主星的 95\%；Patroclus–Menoetius、Altjira 与 Lempo–Hiisi 的尺寸约为 90\%，其中 Lempo–Paha 约为 50\%。已知绝对尺寸最大的小行星卫星是冥王星的卫星 Charon，其直径约为冥王星的一半。

此外，在若干遥远的天体周围，人们还发现了多个环系统（参见：Chariklo 与 Chiron 的环）。
\subsubsection{术语}
除了使用“卫星”和“月”这两个术语之外，“双星”（binary minor planet）一词有时也用于描述拥有一颗卫星的小行星，而“三级系统”（triple）用来描述拥有两颗卫星的小行星。如果其中一个天体明显更大，则称其为主星（primary），另一颗称为伴星（secondary）。

当小行星与其卫星尺寸相近时，有时使用“双小行星”（double asteroid）一词；而“binary”一词通常与两者的相对大小无关。当双小行星系统中的两颗天体大小相似时，小行星中心（MPC）将它们称作“binary companions”，而不会简单地将较小者称为卫星。

一个真正意义上的双星系统实例是 90 号 Antiope 系统，其性质于 2000 年 8 月被确认。非常小的卫星则通常称为 moonlets。

\subsection{发现历程}
在哈勃太空望远镜时代和深空探测器抵达太阳系外部之前，寻找小行星卫星的尝试主要依赖于地基光学观测。例如，1978 年，人们通过恒星掩星观测声称发现 532 号 Herculina 拥有一颗卫星。然而，之后哈勃太空望远镜的高分辨率成像并没有发现任何卫星，目前的共识认为 Herculina 并没有显著卫星。

在随后的几年中，还有其他关于小行星伴星的类似报告。当时，天文学家 Thomas Hamilton 在《Sky \& Telescope》杂志中写信指出，地球上某些看似同时形成的撞击坑（例如魁北克的 Clearwater 湖）可能是由成对、相互引力束缚的小天体撞击形成的。

同样在 1978 年，人们发现了冥王星最大的卫星 Charon；但当时冥王星仍被视为一颗主要行星。

1993 年，伽利略号探测器发现一颗小卫星 Dactyl 正环绕小行星 243 号 Ida 运行，这是首个被确认的小行星卫星。第二个例子于 1998 年在小行星 45 号 Eugenia 周围被发现。2001 年，小行星 617 Patroclus 及其等尺寸伴星 Menoetius 成为木星特洛伊群中首个已知的双小行星系统。继冥王星–Charon 系统之后，第一个被清晰分辨的海王星外双星——1998 WW31——于 2002 年被确认。

2021 年，130 号 Elektra 被发现拥有三颗卫星，使其成为目前唯一已知的四重小行星系统。
\subsubsection{多重系统}
在 2005 年，小行星 87 号 Sylvia 被发现拥有两颗卫星，使其成为首个已知的三级系统（也称三级小行星或三级小行星系统）。随后，人们又在小行星 45 号 Eugenia 周围发现了第二颗卫星。同样在 2005 年，矮行星 Haumea 被发现拥有两颗卫星，使其成为继冥王星之后第二个已知拥有多颗卫星的海王星外天体。此外，小行星 216 号 Kleopatra 与 93 号 Minerva 分别在 2008 年和 2009 年被确认为三级小行星系统。目前已知唯一的四重小行星系统是 130 号 Elektra。自最初的多重小行星系统被发现以来，后续仍不断有新的多重系统被确认。截至 2025 年，小行星中的已知多重系统总数达到 18 个（包括冥王星与 Haumea 系统）。

下表列出了所有多重系统的卫星，从冥王星开始。冥王星在其第一颗卫星于 1978 年被发现时尚未获得小行星编号。当前已知的最高重数系统为冥王星（六重系统）与 130 号 Elektra（四重系统）。
\begin{figure}[ht]
\centering
\includegraphics[width=14.25cm]{./figures/6619e07290fe4aba.png}
\caption{} \label{fig_Minor_1}
\end{figure}
\subsection{普遍性}
关于双天体（binary objects）族群的数据仍然不够完整。除了不可避免的观测偏差（例如：与地球的距离、大小、反照率以及组成部分之间的分离度），其出现频率在不同类别的天体中似乎也不尽相同。在小行星中，估计约有 2\% 拥有卫星。在海王星外天体（TNOs）中，估计约有 11\% 是双星或多重系统，而且多数大型 TNO 至少拥有一颗卫星，包括所有由国际天文学联合会列出的四颗矮行星。

在以下三大类天体中，各已知超过 50 个双星系统：近地小行星、小行星带天体与海王星外天体。这些统计尚未包含仅基于光变曲线推断而未被直接成像确认的双星候选体。

在半长轴小于海王星轨道的半人马小行星中，目前已知两例双星系统。二者均为双环系统，分别位于 2060 Chiron 与 10199 Chariklo 周围，并分别于 1993–2011 年与 2013 年发现。
\subsection{起源}
小行星卫星的起源目前尚不能确定，存在多种假说。其中一种模型认为，小行星卫星源自撞击事件中从主天体上被击落的碎片。而另一类配对可能源自小天体被较大天体的引力所俘获。

由碰撞形成的系统受到角动量的限制，即受到质量与相互分离距离的约束。分离距离较近的双星系统符合该模型（例如冥王星–Charon）。然而，对于分离较远且成分大小相近的双星系统来说，除非在形成过程中损失了大量质量，否则其形成不太可能由碰撞解释。

已知小行星双星系统中，两天体之间的距离从数百公里（如 243 Ida、3749 Balam）到超过 3000 千米（如 379 Huenna）不等。在海王星外天体中，已知的分离距离范围约为 3000 到 50000 千米。
\subsection{族群与类别}
\begin{figure}[ht]
\centering
\includegraphics[width=8cm]{./figures/ceb9ae31e27f5fc2.png}
\caption{} \label{fig_Minor_3}
\end{figure}
所谓“双星系统的典型特征”，往往取决于其在太阳系中的位置（推测原因在于，不同小行星族群的形成机制与系统寿命不同）。
\begin{itemize}
\item 在近地小行星中，卫星通常在距离主天体约 3--7 个主星半径的位置运行，其直径比主星小两倍到数倍不等。由于这些双星系统全都属于内侧行星轨道穿越者，人们认为母体在接近行星时受到的潮汐应力可能促成了其中许多双星系统的形成，尽管碰撞也被认为是一些卫星的成因。
\item 在主带小行星中，卫星通常比主星小得多（著名的例外是 90 号 Antiope），运行距离约为主星半径的 10 倍。这里的许多双星系统都是某些小行星家族的成员，其中相当一部分卫星被认为是母体在小行星碰撞事件中解体后形成的碎片，这些碎片共同构成了主天体与其卫星。
\item 在海王星外天体（TNOs）中，绕行的两个组分通常具有相近的尺寸，其轨道半长轴则往往大得多——约为主星半径的 100 到 1000 倍。预计这些双星系统中有相当一部分属于原始形成的天体。
\item 冥王星已知拥有五颗卫星。其中最大的卫星 Charon，其半径超过冥王星的一半大，并且其轨道足以使其绕行的中心点位于冥王星表面之外。实际上，冥王星与 Charon 都绕其共同质心运行，其中冥王星的轨道完全位于 Charon 的轨道之内；因此，这一系统被非正式地视作“双矮行星”双星系统。冥王星的其他四颗卫星——Nix、Hydra、Kerberos 与 Styx——体积都小得多，绕冥王星–Charon 系统运行。
\item Haumea 有两颗卫星，其半径估计分别约为 155 千米（Hiʻiaka）与 85 千米（Namaka）。
\item Makemake 已知有一颗卫星，记为 S/2015 (136472) 1，其直径估计约为 160 千米（约 100 英里）。
\item 47171 Lempo 是一个独特的海王星外三级系统：Lempo 与其质量近似相等的卫星 Hiisi 构成一个近距离双星系统，两者之间相距约 867 千米。第二颗卫星 Paha 则以约 7411 千米的距离绕 Lempo–Hiisi 双星运行。
\item Eris 已知拥有一颗卫星 Dysnomia。根据其亮度推算，其半径估计大约在 150 到 350 千米之间。
\end{itemize}
\subsection{列表}
截至 2022 年 1 月，已知共有 457 个小行星系统，包含 477 个已确认的伴星。下表列出了这些系统按轨道类别划分的总数量：
\begin{table}[ht]
\centering
\caption{}\label{tab_Minor1}
\begin{tabular}{|c|c|c|c|}
\hline
\textbf{系统数量} & \textbf{轨道类别} & \textbf{按类别列表} & \textbf{多重卫星}\\
\hline
86 & 近地天体 & 查看列表 & 三个拥有两颗卫星的系统：3122 Florence、(136617) 1994 CC 和 (153591) 2001 SN263。\\
\hline
31 & 火星穿越小行星 & 查看列表 & 一个拥有两颗卫星的系统：2577 Litva。\\
\hline
212 & 主带小行星 & 查看列表 & 八个拥有两颗卫星的系统：45 Eugenia、87 Sylvia、93 Minerva、107 Camilla、216 Kleopatra、3749 Balam、4666 Dietz、6186 Zenon；以及一个拥有三颗卫星的四重系统：130 Elektra。\\
\hline
6 & 木星特洛伊小行星 & 查看列表 & –\\
\hline
122 & 海王星外天体 & 查看列表 & 两个拥有两颗卫星的系统：47171 Lempo 和 136108 Haumea；一个拥有五颗卫星的系统：134340 Pluto。\\
\hline
\end{tabular}
\end{table}
\subsubsection{近地天体}
以下为拥有伴星的近地小行星列表。状态尚未确认的双星候选体将以深色背景显示。概览内容可参见摘要与引言部分。
\begin{figure}[ht]
\centering
\includegraphics[width=14.25cm]{./figures/04b1eabfa5b8294a.png}
\caption{} \label{fig_Minor_4}
\end{figure}
\begin{figure}[ht]
\centering
\includegraphics[width=14.25cm]{./figures/36a9021ec75f4363.png}
\caption{} \label{fig_Minor_5}
\end{figure}
\begin{figure}[ht]
\centering
\includegraphics[width=14.25cm]{./figures/2fcf746c45764c58.png}
\caption{} \label{fig_Minor_6}
\end{figure}
\begin{figure}[ht]
\centering
\includegraphics[width=14.25cm]{./figures/426940035afd6fa4.png}
\caption{} \label{fig_Minor_7}
\end{figure}
\begin{figure}[ht]
\centering
\includegraphics[width=14.25cm]{./figures/8520d8298395b51f.png}
\caption{} \label{fig_Minor_8}
\end{figure}
\begin{figure}[ht]
\centering
\includegraphics[width=14.25cm]{./figures/ae2131988176812a.png}
\caption{} \label{fig_Minor_9}
\end{figure}
\begin{figure}[ht]
\centering
\includegraphics[width=14.25cm]{./figures/fc7e1f68f72b9e57.png}
\caption{} \label{fig_Minor_10}
\end{figure}
\begin{figure}[ht]
\centering
\includegraphics[width=14.25cm]{./figures/5239278dce824d78.png}
\caption{} \label{fig_Minor_11}
\end{figure}
\subsubsection{火星穿越天体}
以下为拥有伴星的火星穿越小行星列表。状态尚未确认的双星候选体将以深色背景显示。概览内容可参见摘要与引言部分。
\begin{figure}[ht]
\centering
\includegraphics[width=14.25cm]{./figures/a2023f17144bd91d.png}
\caption{} \label{fig_Minor_12}
\end{figure}
\begin{figure}[ht]
\centering
\includegraphics[width=14.25cm]{./figures/1cbd02266bf93b10.png}
\caption{} \label{fig_Minor_13}
\end{figure}
\begin{figure}[ht]
\centering
\includegraphics[width=14.25cm]{./figures/0f36f996d1086b13.png}
\caption{} \label{fig_Minor_14}
\end{figure}
\begin{figure}[ht]
\centering
\includegraphics[width=14.25cm]{./figures/bd7e36185597f0da.png}
\caption{} \label{fig_Minor_15}
\end{figure}
\subsubsection{主带小行星}
以下为拥有伴星的主带小行星列表。状态尚未确认的双星候选体将以深色背景显示。概览内容可参见摘要与引言部分。
\begin{figure}[ht]
\centering
\includegraphics[width=14.25cm]{./figures/b5a7a182eab99d09.png}
\caption{} \label{fig_Minor_16}
\end{figure}
\begin{figure}[ht]
\centering
\includegraphics[width=14.25cm]{./figures/65a834d257ced968.png}
\caption{} \label{fig_Minor_17}
\end{figure}
\begin{figure}[ht]
\centering
\includegraphics[width=14.25cm]{./figures/8f53999a2a1bc882.png}
\caption{} \label{fig_Minor_18}
\end{figure}
\begin{figure}[ht]
\centering
\includegraphics[width=14.25cm]{./figures/822f23a22af727bf.png}
\caption{} \label{fig_Minor_19}
\end{figure}
\begin{figure}[ht]
\centering
\includegraphics[width=14.25cm]{./figures/cab86d1be41506dc.png}
\caption{} \label{fig_Minor_20}
\end{figure}
\begin{figure}[ht]
\centering
\includegraphics[width=14.25cm]{./figures/1d4fe96a782c4a6f.png}
\caption{} \label{fig_Minor_21}
\end{figure}
\begin{figure}[ht]
\centering
\includegraphics[width=14.25cm]{./figures/fcdeabbdb33d5eba.png}
\caption{} \label{fig_Minor_22}
\end{figure}
\begin{figure}[ht]
\centering
\includegraphics[width=14.25cm]{./figures/31ef1982ccf24447.png}
\caption{} \label{fig_Minor_23}
\end{figure}
\begin{figure}[ht]
\centering
\includegraphics[width=14.25cm]{./figures/e6bf37171f0a9543.png}
\caption{} \label{fig_Minor_24}
\end{figure}
\begin{figure}[ht]
\centering
\includegraphics[width=14.25cm]{./figures/c9ae541fd8ad7a26.png}
\caption{} \label{fig_Minor_25}
\end{figure}
\begin{figure}[ht]
\centering
\includegraphics[width=14.25cm]{./figures/12154bf71f047215.png}
\caption{} \label{fig_Minor_26}
\end{figure}
\begin{figure}[ht]
\centering
\includegraphics[width=14.25cm]{./figures/691c5864b76fa250.png}
\caption{} \label{fig_Minor_27}
\end{figure}
\begin{figure}[ht]
\centering
\includegraphics[width=14.25cm]{./figures/fe020baedce774f4.png}
\caption{} \label{fig_Minor_28}
\end{figure}
\begin{figure}[ht]
\centering
\includegraphics[width=14.25cm]{./figures/49e8c08b4f236073.png}
\caption{} \label{fig_Minor_29}
\end{figure}
\begin{figure}[ht]
\centering
\includegraphics[width=14.25cm]{./figures/48b89c64ff188f6b.png}
\caption{} \label{fig_Minor_30}
\end{figure}
\begin{figure}[ht]
\centering
\includegraphics[width=14.25cm]{./figures/6fd2d3261a3dac72.png}
\caption{} \label{fig_Minor_31}
\end{figure}
\begin{figure}[ht]
\centering
\includegraphics[width=14.25cm]{./figures/973998c412b3f3ae.png}
\caption{} \label{fig_Minor_32}
\end{figure}
\begin{figure}[ht]
\centering
\includegraphics[width=14.25cm]{./figures/2001db74213ab6a3.png}
\caption{} \label{fig_Minor_33}
\end{figure}
\begin{figure}[ht]
\centering
\includegraphics[width=14.25cm]{./figures/dba5464c1ab5b6d7.png}
\caption{} \label{fig_Minor_34}
\end{figure}
\begin{figure}[ht]
\centering
\includegraphics[width=14.25cm]{./figures/9751958feb287f9e.png}
\caption{} \label{fig_Minor_35}
\end{figure}
\begin{figure}[ht]
\centering
\includegraphics[width=14.25cm]{./figures/f1f654c9f15f254d.png}
\caption{} \label{fig_Minor_36}
\end{figure}

以下列出的双星系统属于双小行星（double asteroids），其两个组分尺寸相近，且系统的质心位于较大天体之外。
\begin{enumerate}
\item 90 Antiope – S/2000 (90) 1
\item 854 Frostia – 未指定
\item 1313 Berna – 未指定
\item 2478 Tokai – 未指定
\item 3169 Ostro – 未指定
\item 3749 Balam – S/2002 (3749) 1
\item 3905 Doppler – 未指定
\item 4674 Pauling – S/2004 (4674) 1
\item 4951 Iwamoto – 未指定
\item 5674 Wolff – 未指定
\item 8474 Rettig – 未指定
\item 17246 Christophedumas – S/2004 (17246) 1
\item (300163) 2006 VW139 – 未指定
\end{enumerate}
此外，下列天体可能属于双小行星，但由于其尺寸和轨道存在误差，因此尚不确定。
\begin{enumerate}
\item 809 Lundia – 未指定
\item 1089 Tama – 未指定
\item 1509 Esclangona – S/2003 (1509) 1
\item 4492 Debussy – 未指定
\item 11264 Claudiomaccone – 未指定
\item 22899 Alconrad – (22899) Alconrad I Juliekaibarreto
\end{enumerate}
\subsubsection{木星特洛伊小行星}
以下为拥有伴星的木星特洛伊小行星列表。状态尚未确认的双星候选体将以深色背景显示。概览内容可参见摘要与引言部分。


